\documentclass[10pt,a4paper]{amsart}
\usepackage[margin=19mm]{geometry}
\usepackage{amsmath,amssymb,mathtools}
\usepackage[scr=rsfs,cal=euler]{mathalfa}
\usepackage{xfrac}
\usepackage[unicode,hidelinks,bookmarks=false]{hyperref}
\newcommand{\ie}{i.e.\ }
\newcommand{\cf}{cf.\ }
\newcommand{\resp}{respectively\ }
\newcommand{\Subsection}{\subsection}
\providecommand{\nobreakdash}{\nobreak\hbox{-}}
\setlength{\emergencystretch}{2em}
\multlinegap0pt
\usepackage{mathtools}
\usepackage{xcolor}
\usepackage{algorithm}\usepackage{algpseudocode}
\usepackage{siunitx}
\newtheorem{proposition}{Proposition}
\title{\LARGE \bf Linear-Quadratic Gaussian Games with Distributed Sparse Estimation}
\author{Tianyu Qiu, Filippos Fotiadis, Xinjie Liu, Christian Ellis, Jesse Milzman, Wesley Suttle, Ufuk Topcu, David Fridovich-Keil}
\date{Source: arXiv:2603.17202v1 (CC BY 4.0)}
\begin{document}
\maketitle

\noindent\emph{Source and licence.} \LARGE \bf Linear-Quadratic Gaussian Games with Distributed Sparse Estimation. Version arXiv:2603.17202v1 (CC BY 4.0), distributed by arXiv under the Creative Commons Attribution 4.0 International licence, http://creativecommons.org/licenses/by/4.0/, as recorded in the arXiv licence metadata for this exact version and shown on the abstract page https://arxiv.org/abs/2603.17202v1. Full licence notice: \texttt{licenses/arxiv-2603.17202-CC-BY-4.0.txt}.\par\smallskip

\noindent\textit{Editorial scope.} Counted unit: the article's proposition proposition (source0.tex lines 334--336). The article's complete original proof of it is delivered with it (source0.tex lines 337--352, one \texttt{proof} environment of 16 lines). No mathematics is written, repaired or re-derived: the statement and the proof are the article's own lines, in the article's own notation, and no retained line is substituted or re-flowed.  Retained uncounted as the source-local context the proof needs: lines 327--330.  Nothing was omitted from any retained span, so the package carries no omission record.  The retained lines cite 2 works, whose entries are transcribed verbatim from the article's own bibliography source in the e-print and delivered with short editorial labels recorded in the item's reference mapping.  No retained line contains a figure environment, an image-inclusion command or drawing code, so no image is delivered and the package declares no asset dependency.  Editorial formatting only, in addition: an AMS \texttt{amsart} preamble that restates the article's own declarations and macros used by the retained lines, namely the environments \texttt{proposition} on separate counters, together with the standard packages those lines need.  The retained environments print in this document as Proposition 1 in order of appearance, whereas the article prints the same items with the numbering of its own full document.  The complete native source of the article as distributed in the arXiv e-print for this exact version is bundled unchanged as \texttt{sources/196610/source0.tex}.
\begin{equation}
        \hspace{-1mm}\min_{K_t^i}\ \underbrace{\left\|K_t^i(C_t^i\Sigma_t^{i-}C_t^{i\top}+V_t^i)-\Sigma_t^{i-}C_t^i\right\|_F^2}_{L_{\text{Estimation}}}+\underbrace{\sum_{\rho\in[P_i]}\lambda_t^{i\rho}\left\|K_t^i[\rho]\right\|_F}_{L_{\text{Sparsity}}},
    \label{eqn:lasso_sparse_estimation}
\end{equation}

\begin{proposition}
Problem \eqref{eqn:lasso_sparse_estimation} is equivalent to a convex conic program with quadratic cost.
\end{proposition}

\begin{proof}
Using matrix vectorization \cite{macedo2013typing}, the group lasso problem \eqref{eqn:lasso_sparse_estimation} can be reconstructed as
\begin{multline}
        \min_{K_t^i}\left\|\left((C_t^i\Sigma_t^{i-}C_t^{i\top}+V_t^i)\otimes I\right)\text{vec}(K_t^i)-\text{vec}(\Sigma_t^{i-}C_t^i)\right\|_2^2\\
        +\sum_{\rho\in[P_i]}\lambda_t^{i\rho}\|\text{vec}(K_t^i[\rho])\|_2,\label{eqn:vec_group_lasso}
\end{multline}
where $\otimes$ denotes the Kronecker product. Moreover, introducing the slack variables $s_t^{i\rho}$, it follows that the optimization \eqref{eqn:vec_group_lasso} has the same optimal solution as 
\begin{align*}
        \min_{K_t^i, s_t^i}&\left\|\left((C_t^i\Sigma_t^{i-}C_t^{i\top}{+}V_t^i)\otimes I\right)\text{vec}(K_t^i){-}\text{vec}(\Sigma_t^{i-}C_t^i)\right\|_2^2{+}\sum_{\rho\in[P_i]}\lambda_t^{i\rho}s_t^{i\rho},\label{eqn:second-order_cone_optimization}\\
        \text{s.t.}\ &\begin{bmatrix}
            \text{vec}(K_t^i[\rho]) \\
            s_t^{i\rho}
        \end{bmatrix}\in\mathfrak{C},\ \mathfrak{C}\coloneqq\left\{(\mathfrak{p},\mathfrak{q})\left|\ \mathfrak{q}\in\mathbb{R}, \|\mathfrak{p}\|_2\leq \mathfrak{q}\right\}\right.,
\end{align*}
which is a convex conic program with a second-order cone constraint by definition \cite{boyd2004convex}. 
\end{proof}

\begin{thebibliography}{99}
\bibitem[boyd20]{boyd2004convex} S.~Boyd and L.~Vandenberghe, \emph{Convex optimization}.\hskip 1em plus 0.5em minus 0.4em\relax Cambridge university press, 2004.
\bibitem[macedo]{macedo2013typing} H.~D. Macedo and J.~N. Oliveira, ``Typing linear algebra: A biproduct-oriented approach,'' \emph{Science of Computer Programming}, vol.~78, no.~11, pp. 2160--2191, 2013.
\end{thebibliography}
\end{document}
